\documentclass[11pt]{article}
\usepackage[T1]{fontenc}
\usepackage{lmodern}
\usepackage[margin=1in]{geometry}
\usepackage{longtable}
\usepackage{esoterica}
\usepackage[hidelinks]{hyperref}

% Long typewriter identifiers do not hyphenate; prefer loose lines to overfull ones.
\sloppy

\newcommand\pkg[1]{\textsf{#1}}
\newcommand\cs[1]{\texttt{\textbackslash#1}}
\newcommand\meta[1]{\ensuremath{\langle}\textit{#1}\ensuremath{\rangle}}
\newcommand\oarg[1]{[\meta{#1}]}

\title{The \pkg{esoterica} package\\[0.5ex]
  \large Glyphs for esoteric traditions}
\author{Sigilante}
\date{v0.1.0 \quad 2026-09-13}

\begin{document}
\maketitle

\begin{abstract}
\pkg{esoterica} provides symbols from Freemasonry, the Golden Dawn, Druidry
and related traditions: \MasonicSquareAndCompasses\ \GoldenDawnPentagram\
\DruidAwen\ \DruidTriquetra.  The glyphs are original drawings, not wrappers
around a system font, so they work with pdf\LaTeX, Xe\LaTeX\ and Lua\LaTeX,
scale with the surrounding text and take its colour.
\end{abstract}

\tableofcontents

\section{Quick start}

\begin{verbatim}
\usepackage{esoterica}
...
The \MasonicSquareAndCompasses{} and the \eso{goldendawn/Pentagram}.
\end{verbatim}

\noindent gives: The \MasonicSquareAndCompasses{} and the \eso{goldendawn/Pentagram}.

\section{Loading}

By default every built-in tradition is loaded.  To load only some, name them:

\begin{verbatim}
\usepackage[masonic, druidic]{esoterica}
\usepackage[traditions={masonic, druidic}]{esoterica}   % same thing
\usepackage[none]{esoterica}          % registry only; bring your own glyphs
\usepackage[commands=false]{esoterica} % no \MasonicLevel-style commands
\end{verbatim}

\noindent To load a further tradition later, use
\cs{esotericaload}\{\meta{name}\}.  It reads the file
\texttt{esoterica-\meta{name}.def}.

\section{Using glyphs}

Every glyph has an identifier \meta{tradition}\texttt{/}\meta{Name}.  There are
two ways to set one:

\begin{description}
  \item[\cs{eso}\oarg{options}\{\meta{tradition/Name}\}]
    always available, e.g.\ \verb|\eso{masonic/Level}| \eso{masonic/Level}.
  \item[\cs{\meta{Prefix}\meta{Name}}\oarg{options}]
    generated for each glyph, e.g.\ \verb|\MasonicLevel| \MasonicLevel.
    A command that already exists is never overwritten; the package warns
    and you can fall back to \cs{eso}.
\end{description}

Glyphs work in math mode, where they follow the math style:
$x^{\GoldenDawnFire} \GoldenDawnWater_{\GoldenDawnEarth}$.

\subsection{Options}

Give options to a single glyph in brackets.  To set them for the rest of the
current group, use \cs{esotericasetup}\{\meta{options}\}.

\begin{longtable}{lll}
\bfseries Key & \bfseries Default & \bfseries Meaning\\ \hline
\endhead
\texttt{scale}       & \texttt{1}       & Multiplies size, raise and sidebearing\\
\texttt{size}        & \texttt{0.8em}   & Length of one grid unit (the glyph height)\\
\texttt{raise}       & \texttt{-0.08em} & Vertical shift from the baseline\\
\texttt{sidebearing} & \texttt{0.06em}  & Space on each side of the glyph\\
\texttt{line-width}  & \texttt{0.075}   & Stroke width, in grid units\\
\texttt{color}       & current colour   & Any \pkg{xcolor} colour expression\\
\texttt{tikz}        & empty            & Extra TikZ options for the drawing\\
\texttt{driver}      & \texttt{draw}    & \texttt{draw} or \texttt{unicode} (see below)\\
\texttt{direction}   & \texttt{ltr}     & \texttt{rtl} reverses the text of \cs{esotericascript}\\
\texttt{font}        & current font     & \pkg{fontspec} font for the \texttt{unicode} driver\\
\texttt{cache}       & \texttt{true}    & Reuse drawn boxes\\
\texttt{proof-size}  & \texttt{6em}     & Grid unit used by \cs{esotericaproof}\\
\end{longtable}

\noindent For example,
\begin{verbatim}
\GoldenDawnHexagram[scale=2, color=blue]
{\esotericasetup{line-width=0.11} \MasonicSquareAndCompasses \DruidTriquetra}
\end{verbatim}
\noindent gives \GoldenDawnHexagram[scale=2, color=blue] and bolder strokes
for the whole group:
{\esotericasetup{line-width=0.11}\MasonicSquareAndCompasses\ \DruidTriquetra}.
The TikZ style \texttt{esoterica glyph} is applied to every drawing, so
\verb|\tikzset{esoterica glyph/.style={line cap=butt}}| also works.

\subsection{Drivers}

The \texttt{draw} driver runs the glyph's TikZ drawing and works everywhere.
Where a glyph also has a Unicode codepoint, the \texttt{unicode} driver sets
that character from a font instead, which gives searchable, copyable text.
It needs Xe\LaTeX\ or Lua\LaTeX:

\begin{verbatim}
\usepackage{fontspec}
\esotericasetup{driver=unicode, font=Noto Sans Symbols}
\end{verbatim}

\noindent Each driver falls back to the other: a glyph without a drawing is
set from its codepoint, and a codepoint the font lacks is drawn.

\subsection{Scripts}

Some glyph sets are alphabets or ciphers for whole words.  The command
\begin{quote}
  \cs{esotericascript}\oarg{options}\{\meta{script}\}\{\meta{text}\}
\end{quote}
sets plain text in a script letter by letter.  Letters match regardless of
case, spaces stay spaces, and characters without a glyph are printed
unchanged.  Dee wrote Enochian from right to left; the option
\texttt{direction=rtl} reverses the text.

\begin{verbatim}
\esotericascript{royalarch}{Holiness to the Lord}
\esotericascript[direction=rtl]{enochian}{Madriax}
\end{verbatim}

\noindent gives \esotericascript{royalarch}{Holiness to the Lord} and
\esotericascript[direction=rtl]{enochian}{Madriax}.

\begin{longtable}{ll}
\bfseries Script & \bfseries Key\\ \hline
\endhead
\texttt{royalarch} & \esotericascript{royalarch}{ABCDEFGHIJKLM}\\
                   & \esotericascript{royalarch}{NOPQRSTUVWXYZ}\\
\texttt{enochian}  & \esotericascript{enochian}{ABCDEFGHILM}\\
                   & \esotericascript{enochian}{NOPQRSTUXZ}\\
\end{longtable}

\section{The glyphs}

\esotericatable{masonic}
\esotericatable{goldendawn}
\esotericatable{druidic}

\section{Adding glyphs}

\subsection{Declarations}

A tradition is a namespace with an optional command prefix:

\begin{verbatim}
\DeclareEsotericTradition{martinist}{
  name   = Martinism,
  prefix = Martinist,
}
\end{verbatim}

A glyph belongs to a tradition:

\begin{verbatim}
\DeclareEsotericGlyph{martinist}{Seal}{
  description = {Seal of the Martinist Order},
  width   = 1,        % grid units; the height is always 1
  unicode = 2721,     % optional; hex, with or without U+
  draw    = {
    \draw (0.5,0.5) circle[radius=0.42];
    ...
  },
}
\end{verbatim}

\noindent A script maps letters to glyphs for \cs{esotericascript}:

\begin{verbatim}
\DeclareEsotericScript{royalarch}{
  A = masonic/RoyalArchA,
  B = masonic/RoyalArchB,
  ...
}
\end{verbatim}

\noindent Declarations can appear in a document preamble, in a private file,
or in a module \texttt{esoterica-\meta{name}.def} shipped with the package.
Built-in traditions are listed in \verb|\c__esoterica_modules_clist| in
\texttt{esoterica.sty}.  \cs{IfEsotericGlyphTF}\{\meta{id}\}\{\meta{true}\}\{\meta{false}\}
tests whether a glyph exists.

\subsection{Design conventions}

\begin{itemize}
  \item The drawing grid runs from $x=0$ to the glyph width and from $y=0$ to
    $y=1$.  Only this box counts for spacing; ink outside it overlaps
    neighbouring text.
  \item Keep about half a stroke (0.04) clear of every edge so round caps
    stay inside the box.
  \item Use \verb|\draw| for strokes so every glyph shares the package's
    line width, caps and joins.  Use \verb|\fill| for dots and solid shapes.
  \item For a finer or heavier stroke, scale the package width rather than
    giving an absolute one: \verb|\draw[line width=0.6\pgflinewidth]|.
  \item Do not set colours; glyphs inherit the text colour.
  \item Glyphs that belong together should share geometry.  The four
    elemental triangles all use the frame $y\in[0.06,0.82]$.
\end{itemize}

\subsection{Proofing}

\cs{esotericaproof}\oarg{options}\{\meta{id}\} draws a glyph on its grid with
the box outlined and the baseline in red, and sets a specimen at text size
underneath.  \cs{esotericaproofsheet}\oarg{options}\{\meta{tradition}\} does the
same for every glyph of a tradition.

\begin{center}
  \esotericaproof{masonic/SquareAndCompasses}\qquad
  \esotericaproof{druidic/Triquetra}
\end{center}

\section{Proof sheets}

\subsection{Freemasonry}
\esotericaproofsheet[proof-size=4em]{masonic}
\subsection{Golden Dawn}
\esotericaproofsheet[proof-size=4em]{goldendawn}
\subsection{Druidry}
\esotericaproofsheet[proof-size=4em]{druidic}

\end{document}
